\begin{afterword}
\summaryhead

The essay argues that a belief is worth holding only if it tells you what to expect. Two people arguing about whether an unheard falling tree makes a sound expect exactly the same things, so they are not really disagreeing about the world.

The author grants that good beliefs can be about things we never see, such as atoms or the height of a building. What matters is that they lead to something observable. From gravity and the building's height you can predict when you will hear a dropped ball land. Some beliefs connect only to each other, like a professor's claim that a novelist shows ``alienated resublimation.'' They predict nothing, and the author calls them floating beliefs.

His test is to ask of any belief what it leads you to expect, and better still, what it rules out. If the answer is nothing, give the belief up.

\responsehead

The central idea is sound and useful. ``What would I expect to see if this were true?'' is one of the best questions a person can learn to ask, and the essay makes it memorable. The examples are its strongest part, especially the sound recorder left by the tree and the bowling ball timed against a clock.

The idea is also old, and the essay does not say so. Charles Peirce stated it in 1878: ``Consider what effects, that might conceivably have practical bearings, we conceive the object of our conception to have. Then, our conception of these effects is the whole of our conception of the object.'' In 1907 William James told a story very like the falling tree. A man walks around a tree while a squirrel on the trunk keeps to the far side. Does the man go round the squirrel? ``Which party is right,'' James said, ``depends on what you practically mean by `going round' the squirrel.''

The idea's later history shows where it breaks. In the 1930s the logical positivists turned it into a test of which sentences have meaning, and Karl Popper into a test of which theories count as science. Both ran into the problem W.~V.~O. Quine stated in 1951: ``our statements about the external world face the tribunal of sense experience not individually but only as a corporate body.'' Many beliefs predict nothing on their own and predict a great deal in combination with others. The essay half-sees this with the atoms and the floor, then tells the reader to test ``every question of belief'' one at a time. It also never says what the rule means for mathematics or ethics, whose claims predict no experiences at all.

As writing, it is vivid but uneven. The opening delays its point. The metaphors multiply: map, network, rent, barnacles. The last paragraph says the same thing three times. A teacher would praise the examples, ask the student to engage with the long tradition behind the idea, and ask where the rule stops applying.
\end{afterword}
